\subsection{Propiedades topológicas de los morfismos planos.}
\label{subsection:IV.2.3-fr}

\begin{lemma}[2.3.1]
\EGAReaderTarget{ega4.c04.lem.2-3-1}{IV.2.3.1-fr}
Sean $f:X\to Y$ un morfismo cuasicompacto y cuasiseparado,
$g:Y'\to Y$ un morfismo plano; pongamos $X'=X\times_Y Y'$,
$f'=f_{(Y')}:X'\to Y'$. Entonces, para todo $\mathcal O_X$-módulo
cuasicoherente $\mathscr F$, el homomorfismo canónico
\begin{equation}
\tag{2.3.1.1}
\label{IV.2.3.1.1-fr}
  g^*(f_*(\mathscr F))\to
  f'_*(\mathscr F\otimes_{\mathcal O_Y}\mathcal O_{Y'})
\end{equation}
es biyectivo.
\end{lemma}

Es un caso particular de (\textbf{III}, 1.4.15) (mejorado en (1.7.21)).

\begin{proposition}[2.3.2]
\label{IV.2.3.2-fr}
Sean $S$ un preesquema, $f:X\to Y$ un $S$-morfismo cuasicompacto y
cuasiseparado; sean $Z$ el subpreesquema de $Y$, imagen cerrada de $X$ por
$f$ \textup{((\textbf I, 9.5.3)} y \textup{(1.7.8))}, y $j:Z\to Y$ la inyección canónica, de
modo que $f=j\circ g$, donde $g:X\to Z$ es un morfismo
(loc. cit.). Sea $h:S'\to S$ un morfismo plano, y pongamos
$f'=f_{(S')}:X_{(S')}\to Y_{(S')}$; entonces
$j'=j_{(S')}:Z_{(S')}\to Y_{(S')}$ se identifica con la inyección canónica
del subpreesquema de $Y_{(S')}$, imagen cerrada de $X_{(S')}$ por $f'$.
\end{proposition}

Como el morfismo $Y_{(S')}\to Y$ es plano (2.1.4), es posible limitarse al caso
en que $S=Y$ (\textbf I, 3.3.11); si $f=(\psi,\theta)$, se sabe que $Z$ es el
subpreesquema cerrado de $S$ definido por el ideal (cuasicoherente)
$\mathscr J$ de $\mathcal O_S$, núcleo del homomorfismo
$\theta:\mathcal O_S\to f_*(\mathcal O_X)$ (\textbf I, 9.5.2). Como $h$ es un
morfismo plano, el ideal cuasicoherente $\mathscr J\mathcal O_{S'}$ de
$\mathcal O_{S'}$ se identifica con el núcleo de
$h^*(\theta):\mathcal O_{S'}\to h^*(f_*(\mathcal O_X))$. Ahora bien, si
$f'=(\psi',\theta')$, se comprueba fácilmente (por ejemplo, reduciéndose al
caso en que $Y$ y $Y'$ son afines y utilizando (\textbf I, 2.2.4)) que
$\theta':\mathcal O_{S'}\to f'_*(\mathcal O_{X'})$ es el compuesto del
homomorfismo canónico (2.3.1.1)
$h^*(f_*(\mathcal O_X))\to f'_*(\mathcal O_{X'})$ con $h^*(\theta)$; la
conclusión se deduce, por tanto, de (2.3.1) y de (\textbf I, 9.5.2), puesto que $f'$ es
cuasicompacto y cuasiseparado (1.1.2 y 1.2.2).

\begin{env}[2.3.3]
\label{IV.2.3.3-fr}
Diremos que un morfismo $f:X\to Y$ es \emph{cuasiplano} si existe un
$\mathcal O_X$-módulo cuasicoherente $\mathscr F$ \emph{de tipo finito} que
sea \emph{$f$-plano} y cuyo soporte sea \emph{igual a} $X$. Diremos que $f$ es
\emph{cuasifielmente plano} si es \emph{cuasiplano} y \emph{sobreyectivo}. Todo morfismo
plano (resp. fielmente plano) es cuasiplano (resp. cuasifielmente plano), pues
entonces $\mathscr F=\mathcal O_X$ satisface las condiciones precedentes.
\end{env}

Se deduce inmediatamente de (2.1.4) y de (\textbf I, 9.1.13) que, si $f$ es
cuasiplano, entonces, para todo morfismo $Y'\to Y$, el morfismo
$f_{(Y')}:X\times_Y Y'\to Y'$ es cuasiplano. Asimismo (teniendo en cuenta
(\textbf I, 3.5.2)), si $f$ es cuasifielmente plano, también lo es $f_{(Y')}$.

\begin{proposition}[2.3.4]
\EGAReaderTarget{ega4.c04.prop.2-3-4}{IV.2.3.4-fr}
Sea $f:X\to Y$ un morfismo cuasiplano \textup{(2.3.3)}. Entonces $f$ posee las
propiedades siguientes (que son equivalentes en virtud de \textup{(1.10.4)}):
\begin{enumerate}
  \item[\textup{(i)}] Para todo $x\in X$ y toda generización $y'$ de $y=f(x)$, existe
  una generización $x'$ de $x$ tal que $f(x')=y'$.
  \item[\textup{(ii)}] Para todo $x\in X$, la imagen por $f$ de
  $\operatorname{Spec}(\mathcal O_{X,x})$ es
  $\operatorname{Spec}(\mathcal O_{Y,y})$.
\oldpage[IV-2]{15}
  \item[\textup{(iii)}] Para toda parte cerrada irreducible $Y'$ de $Y$, toda
  componente irreducible de $f^{-1}(Y')$ domina $Y'$.
\end{enumerate}
\end{proposition}

Basta, por ejemplo, demostrar (ii). Por hipótesis, existe un
$\mathcal O_X$-módulo cuasicoherente de tipo finito $\mathscr F$, plano
respecto de $f$ y tal que $\operatorname{Supp}(\mathscr F)=X$. Para todo
$x\in X$, $\mathscr F_x$ es entonces un $\mathcal O_x$-módulo de tipo
finito, no nulo, y además $\mathscr F_x$ es un $\mathcal O_y$-módulo
\emph{plano}, para el homomorfismo
$\rho:\mathcal O_y\to\mathcal O_x$. Como este último es local y
$\mathscr F_x\neq0$, el lema de Nakayama muestra que
$\mathscr F_x\otimes_{\mathcal O_y}\boldsymbol{k}(y)\neq0$, por lo que $\mathscr F_x$
es un $\mathcal O_y$-módulo \emph{fielmente plano}
$(\mathbf{0}_{\mathrm I},6.4.1)$. De ello se sigue que, para todo ideal primo
$\mathfrak q$ de $\mathcal O_y$, existe un ideal primo $\mathfrak p$ de
$\mathcal O_x$ tal que $\mathfrak q=\rho^{-1}(\mathfrak p)$
$(\mathbf{0}_{\mathrm I},6.5.1)$, lo que demuestra (ii).

\begin{corollary}[2.3.5]
\label{IV.2.3.5-fr}
Sea $f:X\to Y$ un morfismo que satisfaga las condiciones equivalentes \textup{(i),
(ii), (iii)} de \textup{(2.3.4) (en particular un morfismo \emph{cuasiplano} o un morfismo
\emph{abierto} (1.10.4))}.
\begin{enumerate}
  \item[\textup{(i)}] Sean $Z$, $Z'$ dos partes cerradas irreducibles de $Y$ tales que
  $Z\subset Z'$, y sea $T$ una componente irreducible de $f^{-1}(Z)$;
  entonces existe una componente irreducible $T'$ de $f^{-1}(Z')$ que contiene
  a $T$ (y domina $Z'$).
  \item[\textup{(ii)}] Para toda componente irreducible $T$ de $X$,
  $\overline{f(T)}$ es una componente irreducible de $Y$.
  \item[\textup{(iii)}] Supongamos que $Y$ es irreducible y, si $y$ es su punto
  genérico, supongamos que $f^{-1}(y)$ es irreducible. Entonces $X$ es irreducible.
\end{enumerate}
\end{corollary}

(i) Basta aplicar (2.3.4, (i)) tomando como $x$ el punto genérico
de $T$ (siendo entonces $y=f(x)$ el punto genérico de $Z$) y como $y'$ el
punto genérico de $Z'$.

(ii) Es claro que $\overline{f(T)}=Z$ es irreducible y, en virtud de (i),
$Z$ no puede estar estrictamente contenida en una parte cerrada irreducible
$Z'$ de $Y$.

(iii) En virtud de (ii), toda componente irreducible de $X$ domina $Y$ y, por
tanto, corta a $f^{-1}(y)$; la conclusión se deduce entonces de
$(\mathbf{0}_{\mathrm I},2.1.8)$.
